% !TEX program = pdflatex
% DS Physique-Chimie — 2nde Bac Pro — Chapitre 10
% Température et capteurs thermiques
% Durée : 30 minutes — Barème : /20
% Notions évaluées :
%   - Conversion Celsius <-> Kelvin (dans les deux sens)
%   - Exploitation d'une courbe d'étalonnage de capteur (CTN, thermocouple)

\documentclass[a4paper,11pt]{article}
\usepackage[utf8]{inputenc}
\usepackage[T1]{fontenc}
\usepackage[french]{babel}
\usepackage{geometry}
\geometry{left=2cm,right=2cm,top=1.8cm,bottom=2cm}
\usepackage{amsmath,amssymb}
\usepackage{enumitem}
\usepackage{array}
\usepackage{booktabs}
\usepackage{xcolor}
\usepackage{tcolorbox}
\usepackage{fancyhdr}
\usepackage{lastpage}
\usepackage{tikz}
\usepackage{pgfplots}
\pgfplotsset{compat=1.17}

\tcbuselibrary{skins,breakable}

% Couleurs
\definecolor{violet}{HTML}{6f42c1}
\definecolor{violetclair}{HTML}{f5f0ff}
\definecolor{orange}{HTML}{d97706}
\definecolor{bleu}{HTML}{0056b3}
\definecolor{vert}{HTML}{2da44e}
\definecolor{rouge}{HTML}{c53030}
\definecolor{gris}{HTML}{718096}

% En-tête / pied
\pagestyle{fancy}
\fancyhf{}
\lhead{\small DS Physique-Chimie — 2nde Bac Pro}
\rhead{\small Ch.\,10 — Température et capteurs}
\cfoot{\thepage\,/\,\pageref{LastPage}}
\renewcommand{\headrulewidth}{0.5pt}

% Boîtes
\newtcolorbox{enonce}{
  colback=violetclair,
  colframe=violet,
  leftrule=4pt,
  rightrule=0pt,toprule=0pt,bottomrule=0pt,
  boxsep=4pt,left=8pt,right=6pt,
  breakable
}
\newtcolorbox{rappel}{
  colback=yellow!6,
  colframe=orange,
  title={\textbf{Rappels utiles}},
  fonttitle=\bfseries,
  breakable
}
\newtcolorbox{consigne}{
  colback=blue!3,
  colframe=bleu,
  title={\textbf{Consignes}},
  fonttitle=\bfseries
}

% Marge pour barème
\newcommand{\pt}[1]{\hfill\textcolor{gris}{\small\textit{/\,#1 pt}}}

\begin{document}

%===============================================================
% EN-TÊTE
%===============================================================
\begin{center}
{\LARGE\bfseries\color{violet} Devoir Surveillé — Physique-Chimie}\\[4pt]
{\large\bfseries Chapitre 10 — Température et capteurs thermiques}\\[6pt]
{\normalsize Seconde Baccalauréat Professionnel — Durée : \textbf{30 minutes} — Barème : \textbf{/20}}
\end{center}

\vspace{2pt}
\noindent
\begin{tabular}{@{}p{0.55\textwidth}p{0.4\textwidth}@{}}
\textbf{Nom :} \dotfill & \textbf{Prénom :} \dotfill \\[6pt]
\textbf{Classe :} \dotfill & \textbf{Date :} \dotfill \\
\end{tabular}

\vspace{4pt}
\begin{consigne}
\begin{itemize}[leftmargin=*,itemsep=1pt,topsep=2pt]
  \item Calculatrice autorisée. Téléphone interdit.
  \item Soigner la présentation, poser les calculs, indiquer les unités.
  \item Les réponses non justifiées ne seront pas prises en compte.
\end{itemize}
\end{consigne}

\begin{rappel}
\begin{itemize}[leftmargin=*,itemsep=1pt,topsep=2pt]
  \item \textbf{Conversion des températures :} $T\,(\text{K}) = \theta\,(^\circ\text{C}) + 273$
        \quad et \quad $\theta\,(^\circ\text{C}) = T\,(\text{K}) - 273$
  \item \textbf{Zéro absolu :} $0$\,K $= -273\ ^\circ$C (température la plus basse possible).
  \item \textbf{Courbe d'étalonnage :} relie la grandeur électrique mesurée (résistance $R$ en $\Omega$, tension $U$ en mV) à la température réelle.
  \item \textbf{Interpolation linéaire} entre deux points $(T_1,y_1)$ et $(T_2,y_2)$ :
        \[ y = y_1 + \frac{T - T_1}{T_2 - T_1}\,(y_2 - y_1) \]
\end{itemize}
\end{rappel}

%===============================================================
% EXERCICE 1 — CONVERSIONS
%===============================================================
\vspace{4pt}
\noindent{\large\bfseries\color{violet}Exercice 1 — Du Celsius au Kelvin et inversement \pt{6}}

\begin{enonce}
Un technicien d'atelier travaille sur une étuve de séchage du bois. Il doit remplir une fiche de maintenance dans laquelle les températures doivent être exprimées en Kelvin, alors que ses instruments affichent des degrés Celsius (et inversement).
\end{enonce}

\begin{enumerate}[label=\textbf{\arabic*.},leftmargin=*,itemsep=4pt]

\item L'étuve de séchage est réglée sur la consigne $\theta = 65\ ^\circ$C. Exprimer cette température en Kelvin. \pt{1}

\vspace{1.2cm}

\item La température ambiante de l'atelier est $\theta = 18\ ^\circ$C. Convertir cette température en Kelvin. \pt{1}

\vspace{1.2cm}

\item Une sonde Pt100 indique $T = 313$\,K à la sortie d'un four. Exprimer cette température en degrés Celsius. \pt{1}

\vspace{1.2cm}

\item Un entrepôt frigorifique de stockage de colles est maintenu à $T = 268$\,K. Convertir cette température en degrés Celsius. \pt{1}

\vspace{1.2cm}

\item Compléter le tableau de correspondance ci-dessous. \pt{1,5}

\vspace{2pt}
\begin{center}
\renewcommand{\arraystretch}{1.4}
\begin{tabular}{|>{\columncolor{violetclair}}c|c|c|c|c|c|}
\hline
\rowcolor{violetclair}
$\theta\ (^\circ$C) & $-273$ & $0$ & \rule{1.3cm}{0.3pt} & $100$ & \rule{1.3cm}{0.3pt} \\
\hline
$T\ $(K) & \rule{1.3cm}{0.3pt} & \rule{1.3cm}{0.3pt} & $293$ & \rule{1.3cm}{0.3pt} & $573$ \\
\hline
\end{tabular}
\end{center}

\vspace{4pt}

\item \textbf{Justifier} pourquoi il est impossible de mesurer une température de $-10$\,K dans l'atelier. \pt{0,5}

\vspace{1.2cm}
\end{enumerate}

\newpage

%===============================================================
% EXERCICE 2 — EXPLOITATION COURBE CTN
%===============================================================
\noindent{\large\bfseries\color{violet}Exercice 2 — Exploitation d'une courbe d'étalonnage d'une CTN \pt{9}}

\begin{enonce}
Dans une étuve de séchage du bois, la température est mesurée par une \textbf{thermistance CTN}. Le tableau et la courbe d'étalonnage ci-dessous donnent la résistance $R$ de la CTN en fonction de la température $\theta$.
\end{enonce}

\begin{center}
\renewcommand{\arraystretch}{1.3}
\begin{tabular}{|>{\columncolor{violetclair}}c|c|c|c|c|c|}
\hline
\rowcolor{violetclair}
$\theta\ (^\circ$C) & $20$ & $40$ & $60$ & $80$ & $100$ \\
\hline
$R\ (\Omega)$ & $10\,000$ & $5\,000$ & $2\,000$ & $1\,000$ & $500$ \\
\hline
\end{tabular}
\end{center}

\vspace{4pt}
\begin{center}
\begin{tikzpicture}
\begin{axis}[
  width=13.5cm, height=7.5cm,
  xlabel={Température $\theta$ (en $^\circ$C)},
  ylabel={Résistance $R$ (en $\Omega$)},
  xmin=10, xmax=110,
  ymin=0, ymax=11000,
  xtick={20,30,40,50,60,70,80,90,100},
  ytick={0,1000,2000,3000,4000,5000,6000,7000,8000,9000,10000},
  grid=both,
  major grid style={line width=.4pt,draw=gray!40},
  minor grid style={line width=.2pt,draw=gray!20},
  title={\textbf{Courbe d'étalonnage de la CTN}},
  every axis plot/.append style={thick},
]
\addplot[violet,mark=*,mark size=2.2pt,mark options={fill=violet}]
  coordinates {(20,10000) (40,5000) (60,2000) (80,1000) (100,500)};
\end{axis}
\end{tikzpicture}
\end{center}

\vspace{2pt}
\begin{enumerate}[label=\textbf{\arabic*.},leftmargin=*,itemsep=4pt]

\item Comment évolue la résistance $R$ de la CTN lorsque la température $\theta$ augmente ? En déduire la signification du sigle \textbf{CTN}. \pt{1,5}

\vspace{1.4cm}

\item L'ohmmètre branché sur la CTN affiche $R = 5\,000\ \Omega$. En utilisant la courbe d'étalonnage, déterminer la température $\theta$ dans l'étuve. \pt{1,5}

\vspace{1.4cm}

\item La consigne de séchage du bois est fixée à $\theta = 80\ ^\circ$C. Quelle est la valeur de la résistance $R$ que doit détecter le régulateur de l'étuve ? \pt{1,5}

\vspace{1.4cm}

\item Convertir cette température de consigne ($80\ ^\circ$C) en Kelvin. \pt{1}

\vspace{1.2cm}

\item On mesure maintenant une résistance $R = 3\,500\ \Omega$. Cette valeur ne figure pas dans le tableau.\par
\textbf{a.} Entre quelles deux valeurs de température se situe $\theta$ ? \pt{1}

\vspace{1.2cm}

\textbf{b.} En utilisant l'\textbf{interpolation linéaire} entre $\theta_1 = 40\ ^\circ$C ($R_1 = 5\,000\ \Omega$) et $\theta_2 = 60\ ^\circ$C ($R_2 = 2\,000\ \Omega$), calculer la valeur de la température mesurée. \pt{2,5}

\vspace{2.2cm}
\end{enumerate}

\newpage
%===============================================================
% EXERCICE 3 — THERMOCOUPLE
%===============================================================
\noindent{\large\bfseries\color{violet}Exercice 3 — Thermocouple type K dans un four industriel \pt{5}}

\begin{enonce}
Dans un four de traitement thermique, la température est mesurée par un \textbf{thermocouple de type K}. La tension $U$ délivrée (en mV) est proportionnelle à la température $\theta$ (en $^\circ$C), selon la relation :
\[ U\ (\text{mV}) = 0{,}04 \times \theta\ (^\circ\text{C}) \]
\end{enonce}

\begin{enumerate}[label=\textbf{\arabic*.},leftmargin=*,itemsep=4pt]

\item Calculer la tension $U$ délivrée par le thermocouple pour $\theta = 250\ ^\circ$C. \pt{1,5}

\vspace{1.6cm}

\item À la sortie du four, le voltmètre affiche $U = 16$\,mV. En déduire la valeur de la température $\theta$ en degrés Celsius. \pt{1,5}

\vspace{1.6cm}

\item Exprimer cette température en Kelvin. \pt{1}

\vspace{1.2cm}

\item Parmi les capteurs suivants : \emph{CTN}, \emph{Pt100}, \emph{thermocouple type K}, lequel est le plus adapté pour mesurer une température élevée de $900\ ^\circ$C dans un four ? Justifier brièvement. \pt{1}

\vspace{1.6cm}
\end{enumerate}

\vspace{4pt}
\hrulefill\\[2pt]
\hfill\textit{\small Fin du sujet — Vérifier ses réponses avant de rendre la copie.}

%===============================================================
% CORRECTION (page séparée)
%===============================================================
\newpage
\begin{center}
{\Large\bfseries\color{violet} Correction — DS Chapitre 10}\\[2pt]
{\small Température et capteurs thermiques}
\end{center}

\vspace{4pt}
\noindent{\bfseries\color{violet}Exercice 1 — Conversions} \hfill\textit{\small(/6 pt)}
\begin{enumerate}[label=\textbf{\arabic*.},leftmargin=*,itemsep=2pt]
\item $T = 65 + 273 = \mathbf{338\ \text{K}}$
\item $T = 18 + 273 = \mathbf{291\ \text{K}}$
\item $\theta = 313 - 273 = \mathbf{40\ ^\circ\text{C}}$
\item $\theta = 268 - 273 = \mathbf{-5\ ^\circ\text{C}}$ (congélation)
\item Tableau complété :
\begin{center}
\renewcommand{\arraystretch}{1.25}
\begin{tabular}{|>{\columncolor{violetclair}}c|c|c|c|c|c|}
\hline
\rowcolor{violetclair}
$\theta\ (^\circ$C) & $-273$ & $0$ & $\mathbf{20}$ & $100$ & $\mathbf{300}$ \\
\hline
$T\ $(K) & $\mathbf{0}$ & $\mathbf{273}$ & $293$ & $\mathbf{373}$ & $573$ \\
\hline
\end{tabular}
\end{center}
\item Le zéro absolu vaut $0$\,K : il est impossible d'atteindre une température négative en Kelvin. $-10$\,K n'a donc pas de sens physique.
\end{enumerate}

\vspace{6pt}
\noindent{\bfseries\color{violet}Exercice 2 — Courbe d'étalonnage de la CTN} \hfill\textit{\small(/9 pt)}
\begin{enumerate}[label=\textbf{\arabic*.},leftmargin=*,itemsep=2pt]
\item Lorsque $\theta$ augmente, $R$ \textbf{diminue} (de $10\,000\ \Omega$ à $500\ \Omega$). D'où le nom \textbf{CTN} : Coefficient de Température \textbf{Négatif}.
\item Lecture graphique ou dans le tableau : $R = 5\,000\ \Omega \Rightarrow \mathbf{\theta = 40\ ^\circ\text{C}}$.
\item Pour $\theta = 80\ ^\circ$C : $\mathbf{R = 1\,000\ \Omega}$.
\item $T = 80 + 273 = \mathbf{353\ \text{K}}$.
\item \textbf{a.} $R = 3\,500\ \Omega$ est compris entre $R_1 = 5\,000\ \Omega$ ($40\ ^\circ$C) et $R_2 = 2\,000\ \Omega$ ($60\ ^\circ$C), donc $\theta$ est entre $\mathbf{40\ ^\circ\text{C}}$ et $\mathbf{60\ ^\circ\text{C}}$.\par
\textbf{b.} Interpolation linéaire :
\[ \theta = \theta_1 + \frac{R - R_1}{R_2 - R_1}(\theta_2 - \theta_1)
  = 40 + \frac{3\,500 - 5\,000}{2\,000 - 5\,000}\times(60 - 40) \]
\[ \theta = 40 + \frac{-1\,500}{-3\,000}\times 20 = 40 + 0{,}5\times 20 = \mathbf{50\ ^\circ\text{C}}. \]
\end{enumerate}

\vspace{6pt}
\noindent{\bfseries\color{violet}Exercice 3 — Thermocouple type K} \hfill\textit{\small(/5 pt)}
\begin{enumerate}[label=\textbf{\arabic*.},leftmargin=*,itemsep=2pt]
\item $U = 0{,}04 \times 250 = \mathbf{10\ \text{mV}}$.
\item $\theta = \dfrac{U}{0{,}04} = \dfrac{16}{0{,}04} = \mathbf{400\ ^\circ\text{C}}$.
\item $T = 400 + 273 = \mathbf{673\ \text{K}}$.
\item Le \textbf{thermocouple type K} : seul capteur de la liste capable de mesurer des températures aussi élevées (plage jusqu'à environ $1\,000\ ^\circ$C). La CTN ($\leq 150\ ^\circ$C) et la Pt100 (précise mais plutôt utilisée jusqu'à $\sim 850\ ^\circ$C) sont moins adaptées ou en limite de plage.
\end{enumerate}

\end{document}
